% ============================================================================
% CONSIDERAÇÕES FINAIS — planejamento de rotas
% Chamado em partes/textuais/conclusao.tex.
% ============================================================================

No planejamento de rotas, a ordem de visita dos pontos foi modelada como o Problema do Caixeiro Viajante com custo igual à distância tridimensional do circuito, que, em um VANT, relaciona-se diretamente com o tempo de voo e o consumo de energia. A comparação entre a ordem aleatória, o algoritmo guloso, o algoritmo genético (AG) e dois métodos exatos (a força bruta e o Held-Karp) mostrou que a obtenção do ótimo global deixa de ser viável com poucas dezenas de pontos. A força bruta foi utilizada somente até 10 pontos, pois com 11 a instância já levou 18~s e, com 20, levaria cerca de 32 mil anos. O Held-Karp, que obtém o mesmo ótimo por programação dinâmica, foi utilizado até 22 pontos, com cerca de 1~min por instância; com 23, ultrapassou 2~min e exigiu 1,7~GB, sendo retirado do comparativo, e a partir de 25 a memória necessária (7,5~GB) excedeu a do computador utilizado. Assim, a partir de 23 pontos, a comparação restringiu-se ao AG, ao guloso e à ordem aleatória.

O algoritmo genético mostrou-se viável em toda a faixa avaliada, de 3 a 30 pontos, com tempo de cálculo entre 0,13 e 0,45~s no computador pessoal, alocação de memória inferior a 0,1~MB e rotas com qualidade média entre 96,5\% e 100\% da referência, tendo obtido o ótimo global em 150 das 200 instâncias em que ele era conhecido. O algoritmo guloso, embora mais rápido, produziu rotas de 6\% a 18\% mais longas em média. Por ter sido validado nas faixas em que o ótimo é conhecido, o AG foi adotado também nas instâncias maiores, nas quais o ótimo não pode ser calculado e suas rotas se mantiveram próximas da melhor rota conhecida. A execução de rotas tridimensionais pelo controlador de voo, com o VANT simulado no STM32, confirmou que a ordem calculada é executada conforme planejada, com erro máximo de 3,1~cm em cada ponto e redução de 21,3\% na distância total da rota do AG em relação à do guloso.

O mesmo AG, executado no Raspberry Pi Zero W, produziu rotas idênticas às do computador pessoal, verificadas por assinatura SHA-256, em um intervalo de 10 a 24~s, tempo considerado aceitável para um cálculo pré-voo. O algoritmo genético constituiu, assim, uma solução viável, adotada e validada no protótipo, abrangendo sua execução no microprocessador embarcado e a integração com o controlador de voo. Essa validação corrobora sua viabilidade operacional.

Como trabalhos futuros para esta frente, recomenda-se reduzir o tempo de cálculo no Raspberry Pi, seja pela diminuição do número de gerações, seja pela adoção de uma implementação de código compilado; aprofundar o estudo da busca local \textit{2-opt} com múltiplas partidas, que, na investigação exploratória, obteve rotas superiores às do AG em cerca de 1/30 do tempo, seja como alternativa exclusiva, seja combinada ao AG por meio de um algoritmo memético; e incorporar à função de custo variáveis e restrições dinâmicas de tempo, energia e obstáculos, características inerentes a operações logísticas e missões reais.
